\documentclass[11pt,a4paper]{bxjsarticle}

\usepackage{amsmath,amssymb,amsthm,mathtools}
\usepackage{bm}
\usepackage{booktabs}
\usepackage{longtable}
\usepackage{array}
\usepackage{multirow}
\usepackage{geometry}
\usepackage{enumitem}
\usepackage{hyperref}
\usepackage{float}

\geometry{margin=25mm}

\title{技術者派遣における案件アサイン問題の問題設定整理ノート}
\author{}
\date{}

\begin{document}

\maketitle

\section{このモデルで扱いたい問題}

本ノートでは，1人の派遣技術者を対象として，待機時に
\begin{itemize}[leftmargin=2em]
    \item 高単価案件を狙う
    \item 低単価案件を狙う
    \item 研修を行う
\end{itemize}
のいずれを選ぶと長期的な利益が最大になるかを考える．

この問題は，
\begin{itemize}[leftmargin=2em]
    \item 短期利益と長期利益のトレードオフを含む
    \item 案件提案は確率的に成功・失敗する
    \item 研修は即時利益を生まないが将来の成功率を改善しうる
    \item 配属中は市場ニーズ変化により相対適合度が低下しうる
\end{itemize}
という性質を持つため，マルコフ決定過程（MDP）として定式化する．

\section{モデルの前提}

\subsection{対象}
\begin{itemize}[leftmargin=2em]
    \item 管理対象は1人の技術者である．
    \item 意思決定者はアサイン担当マネージャーである．
    \item 各期ごとに逐次的な意思決定を行う．
\end{itemize}

\subsection{簡略化の前提}
現実を簡略化するため，以下を仮定する．
\begin{enumerate}[leftmargin=2em,label=(\arabic*)]
    \item 技術者は常に
    \[
        \text{高単価案件稼働中},\quad \text{低単価案件稼働中},\quad \text{待機中}
    \]
    のいずれかの状態にある．
    \item 意思決定は待機状態でのみ行う．
    \item 提案成功確率は技術者の適合度 \(k\) に依存する．
    \item 研修は待機状態で行い，確率的に適合度を上げる．
    \item 配属中は，最新技術への学習機会の不足や市場ニーズ変化により，相対適合度が確率的に低下する．
    \item 高単価案件・低単価案件の継続確率は，現時点の簡略版では定数とする．
\end{enumerate}

\section{状態・行動・報酬・遷移の定義}

\subsection{状態}

時刻 \(t\) における状態を
\[
    s_t = (z_t, k_t)
\]
とする．

\subsubsection*{稼働状態}
\[
    z_t \in \{z_B, z_H, z_L\}
\]
とする．それぞれ
\begin{align*}
    z_B &: \text{Bench（待機）},\\
    z_H &: \text{High（高単価案件に稼働中）},\\
    z_L &: \text{Low（低単価案件に稼働中）}
\end{align*}
を表す．

\subsubsection*{適合度}
\[
    k_t \in \mathcal{K}=\{0,1,2,\dots,K\}
\]
とする．ここで \(k_t\) は案件への総合的適合度を表す離散状態であり，技術力・経験・案件とのマッチ度を簡略化してまとめた指標である．
\(k_t\) が大きいほど，高単価案件・低単価案件への提案成功率が高いとする．

\subsection{行動}

状態に応じた行動集合を
\[
    A(z_B)=\{a_H,a_L,a_T\}, \qquad A(z_H)=A(z_L)=\{a_N\}
\]
と定義する．

各行動の意味は以下の通りである．
\begin{align*}
    a_H &: \text{高単価案件へ提案する},\\
    a_L &: \text{低単価案件へ提案する},\\
    a_T &: \text{研修を行う},\\
    a_N &: \text{何もしない（稼働継続中の自動遷移用ダミー行動）}.
\end{align*}

\paragraph{重要な前提}
意思決定は待機状態 \(z_B\) にあるときのみ行う．
高単価案件稼働中 \(z_H\)，低単価案件稼働中 \(z_L\) では行動選択を行わず，案件継続または終了に応じて自動遷移する．

\subsection{報酬}

報酬は遷移前報酬として
\[
    r(s_t,a_t)=\mathrm{rev}(z_t)-C_w-C_{\mathrm{edu}}\mathbf{1}[a_t=a_T]
\]
と定義する．

ここで，
\begin{itemize}[leftmargin=2em]
    \item \(\mathrm{rev}(z_t)\)：状態 \(z_t\) に応じた売上
    \item \(C_w\)：各期に必ず発生する給与コスト
    \item \(C_{\mathrm{edu}}\)：研修時にのみ発生する教育コスト
    \item \(\mathbf{1}[a_t=a_T]\)：\(a_t=a_T\) のとき1，それ以外で0
\end{itemize}
である．

売上は
\[
    \mathrm{rev}(z_t)=
    \begin{cases}
        0   & (z_t=z_B),\\
        R_H & (z_t=z_H),\\
        R_L & (z_t=z_L)
    \end{cases}
\]
と定義する．

したがって，
\begin{itemize}[leftmargin=2em]
    \item 待機状態では売上は0
    \item 高単価案件・低単価案件に稼働中の期には売上が発生
    \item 提案成功によって配属が決まった場合でも，売上は次期以降の稼働状態に応じて発生
\end{itemize}
と解釈する．

\section{遷移モデル}

遷移確率は
\[
    P(s_{t+1}\mid s_t,a_t)=P((z',k')\mid(z,k),a)
\]
で表す．

\subsection{待機状態 \(z_B\) での遷移}

\subsubsection{高単価案件へ提案 \(a_H\)}

提案成功確率を \(p_H(k)\) とする．
\[
    P((z_H,k)\mid(z_B,k),a_H)=p_H(k)
\]
\[
    P((z_B,k)\mid(z_B,k),a_H)=1-p_H(k)
\]

\subsubsection{低単価案件へ提案 \(a_L\)}

提案成功確率を \(p_L(k)\) とする．
\[
    P((z_L,k)\mid(z_B,k),a_L)=p_L(k)
\]
\[
    P((z_B,k)\mid(z_B,k),a_L)=1-p_L(k)
\]

\subsubsection{研修 \(a_T\)}

研修成功確率を \(q_T(k)\) とする．成功時には適合度が1上がるが，上限 \(K\) を超えない．
\[
    P((z_B,\min(k+1,K))\mid(z_B,k),a_T)=q_T(k)
\]
\[
    P((z_B,k)\mid(z_B,k),a_T)=1-q_T(k)
\]

\subsection{高単価案件稼働中 \(z_H\) での遷移}

高単価案件の継続確率を \(d_H\)，配属中の適合度低下確率を \(\delta\) とする．
遷移の順序は次の通りとする．
\begin{enumerate}[leftmargin=2em,label=(\arabic*)]
    \item 現在の \(k_t\) に基づいて継続判定
    \item 継続していた場合のみ，期末に確率 \(\delta\) で適合度が1低下
\end{enumerate}

したがって，
\[
    P((z_H,k)\mid(z_H,k),a_N)=d_H(1-\delta)
\]
\[
    P((z_H,\max(k-1,0))\mid(z_H,k),a_N)=d_H\delta
\]
\[
    P((z_B,k)\mid(z_H,k),a_N)=1-d_H
\]

\subsection{低単価案件稼働中 \(z_L\) での遷移}

低単価案件の継続確率を \(d_L\) とする．
\[
    P((z_L,k)\mid(z_L,k),a_N)=d_L(1-\delta)
\]
\[
    P((z_L,\max(k-1,0))\mid(z_L,k),a_N)=d_L\delta
\]
\[
    P((z_B,k)\mid(z_L,k),a_N)=1-d_L
\]

\section{K=3 のときの状態集合}

試作版として
\[
    K=3
\]
とする．

このとき状態集合は
\[
    S=\{(z,k)\mid z\in\{z_B,z_H,z_L\},\ k\in\{0,1,2,3\}\}
\]
であり，状態数は
\[
    |S|=3\times 4=12
\]
である．

\subsection{状態一覧表}

\begin{longtable}{ccl}
\toprule
state\_id & 状態 \((z,k)\) & 意味 \\
\midrule
\endfirsthead
\toprule
state\_id & 状態 \((z,k)\) & 意味 \\
\midrule
\endhead
0  & \((z_B,0)\) & 待機・適合度0 \\
1  & \((z_B,1)\) & 待機・適合度1 \\
2  & \((z_B,2)\) & 待機・適合度2 \\
3  & \((z_B,3)\) & 待機・適合度3 \\
4  & \((z_H,0)\) & 高単価案件稼働中・適合度0 \\
5  & \((z_H,1)\) & 高単価案件稼働中・適合度1 \\
6  & \((z_H,2)\) & 高単価案件稼働中・適合度2 \\
7  & \((z_H,3)\) & 高単価案件稼働中・適合度3 \\
8  & \((z_L,0)\) & 低単価案件稼働中・適合度0 \\
9  & \((z_L,1)\) & 低単価案件稼働中・適合度1 \\
10 & \((z_L,2)\) & 低単価案件稼働中・適合度2 \\
11 & \((z_L,3)\) & 低単価案件稼働中・適合度3 \\
\bottomrule
\end{longtable}

\subsection{行動一覧表}

\begin{table}[H]
\centering
\begin{tabular}{ccl}
\toprule
action\_id & 行動 & 意味 \\
\midrule
0 & \(a_H\) & 高単価案件へ提案 \\
1 & \(a_L\) & 低単価案件へ提案 \\
2 & \(a_T\) & 研修を行う \\
3 & \(a_N\) & 何もしない（稼働継続中） \\
\bottomrule
\end{tabular}
\end{table}

\subsection{状態ごとの合法行動}

\begin{table}[H]
\centering
\begin{tabular}{cl}
\toprule
状態 & 合法行動 \\
\midrule
\((z_B,k)\) & \(a_H, a_L, a_T\) \\
\((z_H,k)\) & \(a_N\) \\
\((z_L,k)\) & \(a_N\) \\
\bottomrule
\end{tabular}
\end{table}

\section{K=3 のときの遷移ルール一覧}

\subsection{待機状態 \((z_B,k)\)}

\subsubsection*{\((z_B,0)\)}
\begin{table}[H]
\centering
\begin{tabular}{cccc}
\toprule
行動 & 次状態 & 確率 & 報酬 \\
\midrule
\(a_H\) & \((z_H,0)\) & \(p_H(0)\) & \(-C_w\) \\
\(a_H\) & \((z_B,0)\) & \(1-p_H(0)\) & \(-C_w\) \\
\(a_L\) & \((z_L,0)\) & \(p_L(0)\) & \(-C_w\) \\
\(a_L\) & \((z_B,0)\) & \(1-p_L(0)\) & \(-C_w\) \\
\(a_T\) & \((z_B,1)\) & \(q_T(0)\) & \(-C_w-C_{\mathrm{edu}}\) \\
\(a_T\) & \((z_B,0)\) & \(1-q_T(0)\) & \(-C_w-C_{\mathrm{edu}}\) \\
\bottomrule
\end{tabular}
\end{table}

\subsubsection*{\((z_B,1)\)}
\begin{table}[H]
\centering
\begin{tabular}{cccc}
\toprule
行動 & 次状態 & 確率 & 報酬 \\
\midrule
\(a_H\) & \((z_H,1)\) & \(p_H(1)\) & \(-C_w\) \\
\(a_H\) & \((z_B,1)\) & \(1-p_H(1)\) & \(-C_w\) \\
\(a_L\) & \((z_L,1)\) & \(p_L(1)\) & \(-C_w\) \\
\(a_L\) & \((z_B,1)\) & \(1-p_L(1)\) & \(-C_w\) \\
\(a_T\) & \((z_B,2)\) & \(q_T(1)\) & \(-C_w-C_{\mathrm{edu}}\) \\
\(a_T\) & \((z_B,1)\) & \(1-q_T(1)\) & \(-C_w-C_{\mathrm{edu}}\) \\
\bottomrule
\end{tabular}
\end{table}

\subsubsection*{\((z_B,2)\)}
\begin{table}[H]
\centering
\begin{tabular}{cccc}
\toprule
行動 & 次状態 & 確率 & 報酬 \\
\midrule
\(a_H\) & \((z_H,2)\) & \(p_H(2)\) & \(-C_w\) \\
\(a_H\) & \((z_B,2)\) & \(1-p_H(2)\) & \(-C_w\) \\
\(a_L\) & \((z_L,2)\) & \(p_L(2)\) & \(-C_w\) \\
\(a_L\) & \((z_B,2)\) & \(1-p_L(2)\) & \(-C_w\) \\
\(a_T\) & \((z_B,3)\) & \(q_T(2)\) & \(-C_w-C_{\mathrm{edu}}\) \\
\(a_T\) & \((z_B,2)\) & \(1-q_T(2)\) & \(-C_w-C_{\mathrm{edu}}\) \\
\bottomrule
\end{tabular}
\end{table}

\subsubsection*{\((z_B,3)\)}
\begin{table}[H]
\centering
\begin{tabular}{cccc}
\toprule
行動 & 次状態 & 確率 & 報酬 \\
\midrule
\(a_H\) & \((z_H,3)\) & \(p_H(3)\) & \(-C_w\) \\
\(a_H\) & \((z_B,3)\) & \(1-p_H(3)\) & \(-C_w\) \\
\(a_L\) & \((z_L,3)\) & \(p_L(3)\) & \(-C_w\) \\
\(a_L\) & \((z_B,3)\) & \(1-p_L(3)\) & \(-C_w\) \\
\(a_T\) & \((z_B,3)\) & \(1\) & \(-C_w-C_{\mathrm{edu}}\) \\
\bottomrule
\end{tabular}
\end{table}

\subsection{高単価案件状態 \((z_H,k)\)}

\subsubsection*{\((z_H,0)\)}
\begin{table}[H]
\centering
\begin{tabular}{cccc}
\toprule
行動 & 次状態 & 確率 & 報酬 \\
\midrule
\(a_N\) & \((z_H,0)\) & \(d_H\) & \(R_H-C_w\) \\
\(a_N\) & \((z_B,0)\) & \(1-d_H\) & \(R_H-C_w\) \\
\bottomrule
\end{tabular}
\end{table}

\subsubsection*{\((z_H,1)\)}
\begin{table}[H]
\centering
\begin{tabular}{cccc}
\toprule
行動 & 次状態 & 確率 & 報酬 \\
\midrule
\(a_N\) & \((z_H,1)\) & \(d_H(1-\delta)\) & \(R_H-C_w\) \\
\(a_N\) & \((z_H,0)\) & \(d_H\delta\) & \(R_H-C_w\) \\
\(a_N\) & \((z_B,1)\) & \(1-d_H\) & \(R_H-C_w\) \\
\bottomrule
\end{tabular}
\end{table}

\subsubsection*{\((z_H,2)\)}
\begin{table}[H]
\centering
\begin{tabular}{cccc}
\toprule
行動 & 次状態 & 確率 & 報酬 \\
\midrule
\(a_N\) & \((z_H,2)\) & \(d_H(1-\delta)\) & \(R_H-C_w\) \\
\(a_N\) & \((z_H,1)\) & \(d_H\delta\) & \(R_H-C_w\) \\
\(a_N\) & \((z_B,2)\) & \(1-d_H\) & \(R_H-C_w\) \\
\bottomrule
\end{tabular}
\end{table}

\subsubsection*{\((z_H,3)\)}
\begin{table}[H]
\centering
\begin{tabular}{cccc}
\toprule
行動 & 次状態 & 確率 & 報酬 \\
\midrule
\(a_N\) & \((z_H,3)\) & \(d_H(1-\delta)\) & \(R_H-C_w\) \\
\(a_N\) & \((z_H,2)\) & \(d_H\delta\) & \(R_H-C_w\) \\
\(a_N\) & \((z_B,3)\) & \(1-d_H\) & \(R_H-C_w\) \\
\bottomrule
\end{tabular}
\end{table}

\subsection{低単価案件状態 \((z_L,k)\)}

\subsubsection*{\((z_L,0)\)}
\begin{table}[H]
\centering
\begin{tabular}{cccc}
\toprule
行動 & 次状態 & 確率 & 報酬 \\
\midrule
\(a_N\) & \((z_L,0)\) & \(d_L\) & \(R_L-C_w\) \\
\(a_N\) & \((z_B,0)\) & \(1-d_L\) & \(R_L-C_w\) \\
\bottomrule
\end{tabular}
\end{table}

\subsubsection*{\((z_L,1)\)}
\begin{table}[H]
\centering
\begin{tabular}{cccc}
\toprule
行動 & 次状態 & 確率 & 報酬 \\
\midrule
\(a_N\) & \((z_L,1)\) & \(d_L(1-\delta)\) & \(R_L-C_w\) \\
\(a_N\) & \((z_L,0)\) & \(d_L\delta\) & \(R_L-C_w\) \\
\(a_N\) & \((z_B,1)\) & \(1-d_L\) & \(R_L-C_w\) \\
\bottomrule
\end{tabular}
\end{table}

\subsubsection*{\((z_L,2)\)}
\begin{table}[H]
\centering
\begin{tabular}{cccc}
\toprule
行動 & 次状態 & 確率 & 報酬 \\
\midrule
\(a_N\) & \((z_L,2)\) & \(d_L(1-\delta)\) & \(R_L-C_w\) \\
\(a_N\) & \((z_L,1)\) & \(d_L\delta\) & \(R_L-C_w\) \\
\(a_N\) & \((z_B,2)\) & \(1-d_L\) & \(R_L-C_w\) \\
\bottomrule
\end{tabular}
\end{table}

\subsubsection*{\((z_L,3)\)}
\begin{table}[H]
\centering
\begin{tabular}{cccc}
\toprule
行動 & 次状態 & 確率 & 報酬 \\
\midrule
\(a_N\) & \((z_L,3)\) & \(d_L(1-\delta)\) & \(R_L-C_w\) \\
\(a_N\) & \((z_L,2)\) & \(d_L\delta\) & \(R_L-C_w\) \\
\(a_N\) & \((z_B,3)\) & \(1-d_L\) & \(R_L-C_w\) \\
\bottomrule
\end{tabular}
\end{table}

\section{目的関数}

方策 \(\pi\) の下での期待割引累積報酬を
\[
    J^\pi(s)
    =
    \mathbb{E}^{\pi}\left[
        \sum_{t=0}^{\infty}\gamma^t r(s_t,a_t)
        \,\middle|\,
        s_0=s
    \right]
\]
とする．ここで \(\gamma\in[0,1)\) は割引率である．

最適方策は
\[
    \pi^*=\arg\max_\pi J^\pi(s)
\]
で与えられる．

\section{ベースライン方策}

\subsection{ベースライン1：目先稼働率最大化方策}
各待機状態 \((z_B,k)\) において，次期に稼働状態へ移行する確率を最大化する行動を選ぶ．

\subsection{ベースライン2：状態悪化回避方策}
ベースライン1を基本とするが，適合度が十分低い場合は研修を優先する．
例えばある閾値 \(k_{\mathrm{safe}}\) を定め，
\begin{itemize}[leftmargin=2em]
    \item \(k \le k_{\mathrm{safe}}\) なら \(a_T\)
    \item それ以外はベースライン1
\end{itemize}
とする．

\subsection{ベースライン3：決定論的閾値方策}
2つの閾値 \(k_H,k_L\) を用いて
\[
    \pi_{\mathrm{th}}(z_B,k)=
    \begin{cases}
        a_H & (k \ge k_H),\\
        a_L & (k_L \le k < k_H),\\
        a_T & (k < k_L)
    \end{cases}
\]
とする．

\section{評価指標}

100ステップのシミュレーション平均を想定する．

\begin{enumerate}[leftmargin=2em,label=(\arabic*)]
    \item \textbf{平均累積報酬}：100ステップまでの累積報酬の平均．
    \item \textbf{平均稼働率}：100ステップ中，\(z_t\in\{z_H,z_L\}\) である割合．
    \item \textbf{高単価案件状態滞在割合}：100ステップ中，\(z_t=z_H\) である割合．
    \item \textbf{低単価案件状態滞在割合}：100ステップ中，\(z_t=z_L\) である割合．
    \item \textbf{低適合度滞在割合}：100ステップ中，
    \[
        k_t \le \frac{K}{3}
    \]
    を満たす状態にいた割合．適合度空間を大まかに3分割し，その下位領域への滞在時間を見る．
    \item \textbf{平均適合度}：100ステップ中の \(k_t\) の平均値．
\end{enumerate}

\section{実験で振るパラメータ}

感度分析では以下を変化させる．
\begin{itemize}[leftmargin=2em]
    \item \(p_H(k)\)：高単価案件提案成功確率
    \item \(p_L(k)\)：低単価案件提案成功確率
    \item \(q_T(k)\)：研修成功確率
    \item \(\delta\)：配属中の適合度低下確率
    \item \(d_H,d_L\)：案件継続確率
    \item \(R_H-R_L\)：高低単価案件の報酬差
    \item \(\gamma\)：割引率
\end{itemize}

\section{今後の拡張候補}

現時点では大枠は固まっているが，今後必要に応じて以下の拡張が考えられる．

\begin{enumerate}[leftmargin=2em,label=(\arabic*)]
    \item 継続確率を \(k\) 依存にする：
    \[
        d_H(k),\quad d_L(k)
    \]
    \item 高単価案件と低単価案件で適合度低下確率を分ける：
    \[
        \delta_H,\quad \delta_L
    \]
    \item 提案成功時の初期オンボーディングコストや初期報酬を追加する．
\end{enumerate}

\section{まとめ}

本モデルは，待機中の技術者に対して
\[
    \text{高単価案件狙い},\quad \text{低単価案件狙い},\quad \text{研修}
\]
の三択をどう切り替えるべきかを，適合度と案件継続を含む単純なMDPとして表したものである．

構造は以下の通りである．
\begin{itemize}[leftmargin=2em]
    \item 待機時のみ意思決定
    \item 提案成功は確率的
    \item 研修で適合度改善
    \item 配属中は適合度劣化
    \item 稼働中のみ売上発生
\end{itemize}

この形で，価値反復・方策反復・ベースライン比較・Colab実装にそのまま接続できる．

\end{document}
